% Figure: radiation pattern of a small loop antenna, drawn as a
% figure-of-eight in a plane containing the loop axis. The pattern magnitude
% is |cos(t)| measured so the two lobes point left and right (in the plane
% of the loop) and the deep null points up and down (along the loop axis).
% Plain tikz plot of a parametric curve, scaled to a 2.6 cm radius; no
% pgfplots axis, to keep the parametric drawing simple and robust.
\begin{figure}[htbp]
\centering
\begin{tikzpicture}[scale=1.0, >=latex]
  % reference axes through the origin
  \draw[dashed, enggray] (0,-2.9) -- (0,2.9);
  \draw[enggray, thin] (-2.9,0) -- (2.9,0);
  \node[enggray, font=\scriptsize, anchor=south] at (0,2.9)
    {loop axis (null)};
  % figure-of-eight: r(t) = |cos t|, radius scaled to 2.6 cm.
  \draw[engblue, very thick, samples=181, domain=0:360, smooth, variable=\t]
    plot ({2.6*abs(cos(\t))*cos(\t)}, {2.6*abs(cos(\t))*sin(\t)});
  % peak labels in the plane of the loop
  \node[engblue, font=\scriptsize, anchor=west] at (2.62,0.28) {peak};
  \node[engblue, font=\scriptsize, anchor=east] at (-2.62,0.28) {peak};
  % the loop seen edge-on at the origin
  \draw[engred, very thick] (-0.3,0) -- (0.3,0);
  \node[engred, font=\scriptsize, anchor=north] at (0,-0.12)
    {loop (edge-on)};
\end{tikzpicture}
\caption[Radiation pattern of a small loop antenna]{The figure-of-eight
pattern of a small loop antenna in a plane containing the loop axis. The
pattern peaks in the plane of the loop and falls to a sharp null along the
axis, the opposite of the intuition that a loop radiates through its face.
The null is deep and narrow because it comes from a first-order
cancellation of the fields from opposite sides of the loop rather than
from a broad lobe roll-off, which is what makes the small loop the classic
direction-finding antenna: an operator rotates it for a signal null, a
sharper and more repeatable bearing than any peak. The same pattern, read
as a receiving aperture, is why a portable radio's ferrite loopstick must
be turned broadside to a station to hear it best.}
\label{fig:vol04ch10:loop-pattern}
\end{figure}
